%
%
%
%
%
\section*{Data and code availability}\addcontentsline{toc}{section}{Data and code availability}

Every number of this \ifimband volume\else part\fi\ that we computed is generated from a result file of a data deposit or from a cited constant; no computed number is
typed by hand. The data deposit is archived with the collected volume, DOI \href{https://doi.org/10.5281/zenodo.23127547}{10.5281/zenodo.23127547}.

\paragraph{What the deposit contains.} The deposit has two parts. The core holds what the text cites: (i) the result files, in JSON and
plain text, behind every computed number, together with the raw output of the scripts that wrote them; (ii) the scripts, written in Python,
each stating before it runs what it expects to find and the positive and negative controls that it runs; (iii) a manifest with the SHA-256 hash
and the size of every file, which the build checks before it reads a result file; (iv) the decimal expansions of the building blocks of
Part~I that the searches left open (Proposition~\partref{I}{prop:raenge}), including those that the algebraic split of
Remark~\partref{I}{rem:zweiprimitive} settles; (v) the certificates and the second routes for the computations that settle a building block; and
(vi) the file \texttt{OPEN\_PROBLEMS.md}, which lists the open problems of the volume. The archive holds the remaining scripts and results of the
investigation, including negative results that the text does not print.

\paragraph{Where the files are.} The collected volume and its three parts are published as four linked records on Zenodo, and the same files
are in one repository on GitHub. The record of the collected volume (DOI \href{https://doi.org/10.5281/zenodo.23127547}{10.5281/zenodo.23127547}) holds the volume and the whole deposit,
the core and the archive described above. The record of each part holds that part and the result files that its text cites, and links to the
record of the volume. The repository (\href{https://github.com/berkay-yueksel-sayim/erdos-364-365-powerful-numbers-pell}{github.com/berkay-yueksel-sayim/erdos-364-365-powerful-numbers-pell}) holds the four PDFs and the deposit, and its README lists the four records. The records of Parts~I, II and~III have the DOIs \href{https://doi.org/10.5281/zenodo.23127703}{10.5281/zenodo.23127703}, \href{https://doi.org/10.5281/zenodo.23127705}{10.5281/zenodo.23127705} and~\href{https://doi.org/10.5281/zenodo.23127707}{10.5281/zenodo.23127707}.

\paragraph{How the numbers get into the text.} The source of the text refers to each computed number by a key. A build tool reads the result
file named for the key, verifies it against the manifest, computes the number in its printed form, and writes it into the text; it stops if a key
has no source or if a file differs from the manifest. The same tool generates the table of the types of the statements and the list of open
problems from the source of the text, and it is part of the core.

\paragraph{Software.} The scripts use Python~3.12 with gmpy2~2.3, SymPy~1.14 and NumPy~2.4. The elliptic curve method and its relatives are run
with GMP-ECM~7.0.6 \cite{GMPECM706}, used as an installed program inside a container image built from the Debian package and not part of the
deposit; the recipe of the image is in the core. The elliptic curve certificates for primality were produced in the same way with
PARI/GP~2.17.2 \cite{PARI2172} and are part of the deposit; a separate program of ours checks each of them, so the result does not
depend on that program. Further work is planned to support these certificates with a second prover written independently by us. The
text is set with pdf\TeX.

\paragraph{Licenses.} The text and the data are licensed under the Creative Commons Attribution 4.0 International License (CC BY 4.0), and
the code under the Apache License, Version 2.0. Programs of others that the scripts use as tools, such as GMP-ECM, are not part of the
deposit and keep their own licenses.
